\chapter*{Communicating Results from Molecular Docking}

The results from molecular docking and other computational studies aim to inspire new experimental research or provide more insights to biological observations. While experimentalists may not be experts in computational methods, they are well-versed in the biological aspects being investigated. It is essential to communicate findings without using technical jargon that could hinder understanding. It is also important to acknowledge that experimentalists are not entirely unfamiliar with computational techniques. When presenting the results, each score should be explained as clearly as possible with its physiological significance. The experimentalists are more accustomed to reading new results with respect to standard data, in the context of molecular docking, the scores for new molecular designs can be explained in comparison to known drugs or reference molecules. This presents the results with the correct context and self-explanatory manner. 

